% ============================================================
%  Kafedra ilmiy seminari --- taqdimot uchun nutq matni (2-tahrir)
%  Mavzu: O‘zbek tili gaplarini neyron tarmoqlar asosida sintaktik
%         tahlil qilishning matematik modeli
%  Ixtisoslik: 05.01.07 -- Matematik modellashtirish. Sonli usullar
%              va dasturlar majmui
%  Ma’ruzachi: Matlatipov Sanatbek
%  Davomiyligi: 15 daqiqa
%  Slaydlar: Matlatipov_seminar_taqdimot.pdf (22 ta slayd + 1 zaxira slayd)
%
%  2-tahrir: kafedra seminari fikrlari asosida (graflar nazariyasi
%  atamalari, «o‘zak», korpus va daraxt misoli, semantik ahamiyat,
%  o‘zgartirilgan model, n^(n-1) -> O(n^2), Adam bo‘yicha aniq ishlar,
%  «ilk bor»).
%
%  Tamoyil: formulalar va raqamlar slaydda ko‘rsatiladi, ularni
%  to‘liq o‘qish shart emas. Matn qisqa, sodda gaplardan iborat.
%
%  Kompilyatsiya:  pdflatex Matlatipov_seminar_nutq_matni.tex (ikki marta)
% ============================================================
\documentclass[14pt,a4paper]{extarticle}

\usepackage[utf8]{inputenc}
\usepackage[T1]{fontenc}
\usepackage{lmodern}
\usepackage[margin=2.0cm,left=3.0cm,right=2.0cm,top=2.0cm,bottom=2.2cm]{geometry}
\usepackage{setspace}
\usepackage{enumitem}
\usepackage{titlesec}
\usepackage{fancyhdr}
\usepackage{microtype}
\usepackage{parskip}
\usepackage{amsmath,amssymb}
\usepackage{booktabs}
\usepackage{needspace}
\usepackage{array}
\DeclareUnicodeCharacter{02BB}{\textquoteleft}

% slayd sarlavhasi sahifa oxirida yolg‘iz qolmasin
\let\origsection\section
\renewcommand{\section}{\Needspace*{9\baselineskip}\origsection}

\setstretch{1.6}
% o‘/g‘ belgili so‘zlar bo‘g‘inga ajratilmaydi: qatorlar hoshiyaga chiqmasligi uchun
\setlength{\emergencystretch}{3em}
\tolerance=2000
\setlist[itemize]{leftmargin=1.4em,topsep=3pt,itemsep=4pt}
\setlist[enumerate]{leftmargin=1.6em,topsep=3pt,itemsep=3pt}

\titleformat{\section}
   {\normalfont\LARGE\bfseries\filright}
   {}{0em}{}
\titlespacing*{\section}{0pt}{1.4em}{0.3em}

\pagestyle{fancy}
\fancyhf{}
\fancyhead[L]{\small Kafedra ilmiy seminari --- nutq matni (15 daqiqa, 2-tahrir)}
\fancyhead[R]{\small S.~Matlatipov}
\fancyfoot[C]{\thepage}
\renewcommand{\headrulewidth}{0.4pt}
\setlength{\headheight}{17pt}

% [boshlanish -> tugash] jamlanma vaqt, har bir slayd sarlavhasi ostida
\newcommand{\timing}[2]{%
   \par\nobreak\noindent\textbf{[#1\,$\to$\,#2]}\par\vspace{0.3em}}
% urg‘u bilan aytiladigan ibora
\newcommand{\say}[1]{\textbf{#1}}
\newcommand{\pausehere}{\par\vspace{0.4em}\noindent\textbf{[pauza]}\par}
\newcommand{\pointto}[1]{\par\vspace{0.3em}\noindent\textit{\small
   [Ko‘rsating: #1]}\par}

\begin{document}

% ============================================================
%  TAYYORGARLIK SAHIFALARI
% ============================================================
\begin{center}
  {\LARGE\bfseries Nutq matni (2-tahrir)}\\[0.3em]
  {\large «O‘zbek tili gaplarini neyron tarmoqlar asosida sintaktik tahlil
  qilishning matematik modeli»}\\[0.3em]
  {\normalsize Ma’ruzachi: S.~Matlatipov\\
  22 ta slayd (+1 zaxira) \quad$\cdot$\quad 15 daqiqa}
\end{center}

\begin{spacing}{1.12}
\noindent\textbf{Qanday o‘qiladi}
\begin{itemize}[itemsep=1pt]
  \item Matn qisqa, sodda gaplardan iborat --- shoshilmasdan o‘qing (daqiqasiga
        taxminan 95 so‘z). Vaqtlar shu tezlik bo‘yicha hisoblangan.
  \item Formulalar va raqamlar slaydda ko‘rinib turadi; ularni to‘liq o‘qish shart emas.
  \item \textbf{qalin matn} --- urg‘u bilan aytiladi;
        \textbf{[pauza]} --- ikki-uch soniya to‘xtang;
        {\small\itshape [Ko‘rsating: \dots]} --- ekrandagi joyga ishora qiling.
  \item \textbf{«Men»} deb gapiring: korpusni, modelning o‘zgarishini, tajribalarni va
        o‘lchovlarni siz bajargansiz. Algoritm va usul \textbf{nomlarini} (Chu--Liu--Edmonds,
        Adam) aytavering --- bu iqtibos emas, atama («Seminardan oldin tekshiring», 1-band).
  \item Vaqt yetmasa, 5-, 7- va 16-slaydlarni bir-ikki gapga qisqartiring.
\end{itemize}

\noindent\textbf{Vaqt nazorati} \ {\small (har bir bo‘limga shu vaqtda yetib keling)}\\[2pt]
\begin{tabular}{@{}lll@{}}
  \toprule
  Bo‘lim & Slaydlar & Boshlanish\\
  \midrule
  Kirish, masalaning qo‘yilishi, ilmiy yangilik & 1--8 & 0:00\\
  Matematik model ($\Pi$, yoy vaznlari) & 9--11 & 5:38\\
  Sonli usullar (MST, vaqt, Adam) & 12--15 & 8:04\\
  Tajribalar va natijalar & 16--18 & 11:59\\
  Dasturiy majmua, xulosa, aprobatsiya & 19--22 & 13:22\\
  \bottomrule
\end{tabular}

\newpage
\noindent\textbf{Seminar fikrlari va ularning hal qilinishi}\\[2pt]
{\small
\begin{tabular}{@{}>{\raggedright\arraybackslash}p{5.1cm}>{\raggedright\arraybackslash}p{7.9cm}>{\raggedright\arraybackslash}p{1.9cm}@{}}
  \toprule
  \textbf{Fikr} & \textbf{Nima qilindi} & \textbf{Slayd}\\
  \midrule
  «Hokim» noaniq --- graflar nazariyasidagi «o‘zak»; graflar nazariyasi atamalari
    & Hamma joyda: uch, yoy, vazn, kirish darajasi, ildiz, sikl, (maksimal vaznli) qamrovchi
      daraxt, o‘zak va tobe uch. Morfologiyada «o‘zak» o‘rniga «negiz» (adashmaslik uchun)
    & 2, 4, 6, 9--14\\
  Boshida korpus tasviri, daraxt shaklidagi natija va semantik ahamiyat, misol bilan
    & Yangi 2-slayd: korpusdagi haqiqiy CoNLL-U yozuvi $\to$ daraxt. Yangi 3-slayd: «kim? kimga?
      nimani?», «eski kitob do‘koni», erkin so‘z tartibi
    & 2, 3\\
  Sonli usul (Adam): nimaga erishdim, qanday qildim
    & Adam formulalari va qiymatlari; «Men bajargan ishlar» (4 band); natija: yo‘qotish
      $14{,}29\to5{,}35$, LAS $24{,}7\to58{,}2\,\%$ (1\,000 qadam, 375 s); o‘qitish egri chizig‘i
    & 15\\
  Qaysi modelni o‘zgartirdim; $n^{n-1}$ dan $n^2$ ga
    & Stanza bi-affin tahlilchisi, oldin/keyin jadvali ($\mathbb R^{500}\to\mathbb R^{1143}$,
      24,5 $\to$ 27,6 mln parametr); MST slaydi va o‘lchangan vaqt tajribasi
    & 9, 13, 14\\
  Optimallashtirish vaqt bo‘yichami yoki nima bo‘yicha?
    & Ikki masala jadvali: (P1) daraxtni topish --- vaqt bo‘yicha; (P2) o‘qitish --- aniqlik bo‘yicha
    & 12\\
  «Ilk bor» ma’nosida gapirish
    & Ilmiy yangilik va xulosa «ilk bor» bilan (faqat tekshirilsa tasdiqlanadigan joylarda)
    & 8, 20\\
  \bottomrule
\end{tabular}}

\vspace{0.8em}
\noindent\textbf{Mavzu nomi bo‘yicha muqobil variantlar} \ {\small (ilmiy rahbar bilan kelishing)}
\begin{enumerate}[itemsep=1pt]
  \item \textit{(slaydlarda)} «O‘zbek tili gaplarini neyron tarmoqlar asosida sintaktik tahlil
        qilishning matematik modeli».
  \item «Kam resursli o‘zbek tili uchun morfosintaktik tahlilning graf-ehtimollik modeli».
  \item «O‘zbek tili gaplarini bog‘liqlik sintaksisi asosida tahlil qilishning matematik
        modellari, algoritmlari va dasturiy majmuasi».
\end{enumerate}
{\small «Semantik» so‘zini nomga kiritmang: semantik qatlam annotatsiya qilinmagan va baholanmagan
(semantik ahamiyat 3-slaydda misol bilan ko‘rsatiladi).}

\newpage
\noindent\textbf{Savollar uchun asosiy raqamlar} \ {\small (ma’ruzada aytish shart emas)}\\[2pt]
{\small
\begin{tabular}{@{}>{\bfseries\raggedright\arraybackslash}p{3.6cm}>{\raggedright\arraybackslash}p{11.9cm}@{}}
  \toprule
  UzUDT & 684 gap, 7\,582 token; 6 annotator; Cohen $\kappa$: lemma 95,2\,\%, UPOS 93,7\,\%,
          belgilar 90,8\,\%; UD v2.17, validatsiyadan o‘tgan\\
  UT & 500 gap, 5\,930 token; UD v2.15 (2024); UD validatsiyasi: 92 ogohlantirish\\
  Bo‘linmalar & UzUDT tajribalari (.1): 435\,/\,48\,/\,198 gap; birlashma (.2): 781\,/\,78\,/\,325 gap\\
  Korpus statistikasi & noproyektiv daraxtlar 9,2\,\% (UzUDT 11,0\,\%, UT 6,8\,\%); o‘zak uch o‘ngda 84,6\,\%;
          o‘rtacha yoy uzunligi 2,72; qo‘shni so‘zlar orasidagi yoylar 52\,\%; bog‘lanish turlari 39;
          o‘rtacha gap uzunligi 11,4 so‘z\\
  Model & TahrirchiBERT: 12 qatlam, 768, BPE lug‘ati 30\,528; $\mathbf v_i\in\mathbb R^{1143}$
          (FastText variantida $\mathbb R^{500}$); BiLSTM $3\times(2\times400)$; bi-affin 400; teglovchi $2\times(2\times200)$\\
  Parametrlar & Stanza, FastText varianti: 24\,513\,613; o‘zgartirilgan (BERT kanali): 27\,563\,156;
          TahrirchiBERT: 109\,486\,848 (muzlatilgan); jami $\approx$137 mln, o‘qitiladigani 20\,\%\\
  O‘qitish & Adam $(0{,}9;\,0{,}95)$, $\alpha=3\cdot10^{-3}$, $\varepsilon=10^{-6}$, $\|g\|\le1$,
          dropout 0,5, partiya $\le$5\,000 so‘z (UzUDT: 1 davr = 2 qadam; birlashma: 3 qadam);
          3\,000 qadam sabr $\to$ AMSGrad, yana 3\,000 $\to$ to‘xtash\\
  O‘qitish jurnali (E2.1) & eng yaxshi: 1\,000-qadam, 375 s (LAS 58,23, yo‘qotish 5,35); AMSGrad: 4\,000-qadam;
          to‘xtash: 7\,000-qadam, 2\,689 s; yo‘qotish 14,29 $\to$ 4,14 (min 3,73); LAS 7\,000-qadamda 54,19\\
  Eng yaxshi qadam & E1.1: 1\,300; E2.1: 1\,000; E3.1: 1\,000; E1.2: 800; E2.2: 900\\
  Vaqt tajribasi & Apple M2 Max, NumPy, Stanza CLE dasturi. To‘liq saralash: $n=5$ --- 1,45 ms, $n=8$ --- 8,86 s;
          0,52 mks/variant $\Rightarrow$ $n=11$ --- 1,71 sutka, $n=19$ --- $3{,}2\cdot10^{10}$ yil,
          $n=20$ --- $1{,}7\cdot10^{12}$ yil. CLE: $n=11$ --- 0,42 ms, $n=20$ --- 0,93 ms, $n=37$ --- 2,19 ms,
          $n=100$ --- 10,6 ms; o‘sish $n^{1,56}$; test to‘plami (198 gap) --- 0,093 s; to‘g‘rilik 1\,560/1\,560;
          ochko‘z qadamdagi sikllar o‘rtacha: $n=11$ --- 0,93, $n=37$ --- 1,46, $n=100$ --- 2,16\\
  Validatsiya (LAS) & E1.1 55,28; E2.1 \textbf{58,23}; E3.1 55,90; E1.2 62,87; E2.2 \textbf{63,81}\\
  Validatsiya (UPOS) & E1.1 79,19; E2.1 83,54; E3.1 83,39; E1.2 83,78; E2.2 87,19\\
  Bootstrap & oxirgi--o‘rtacha: $\Delta$LAS $=+2{,}64$, 95\,\% IO $[0{,}97;\,4{,}12]$, $p\approx0{,}002$\\
  Test (198 gap) & Stanza E2.1: UAS 71,06, LAS 54,07; spaCy S1.1: UAS 67,72, LAS 45,35\\
  G‘alayonlar ($\Delta$LAS) & diakritika $-1{,}14$; kod almashuvi $-0{,}14$; bosh harf $-2{,}52$;
          tinish belgisiz $-8{,}30$\\
  \bottomrule
\end{tabular}}

\vspace{0.6em}
\noindent\textbf{Seminardan oldin tekshiring} \ {\small (repozitoriy va jurnallar tahlili natijalari)}
\begin{enumerate}[itemsep=2pt]\raggedright
  \item \textbf{Muhim: muallif hissasi qanday aytiladi.} Chu--Liu--Edmonds algoritmi (1965--1967) va
        Adam usuli --- klassik natijalar. Slayd va nutqda «men bu algoritmni yaratdim» deyilmaydi.
        Sizning hissangiz: o‘zbek tili tahlilini graf masalasi ko‘rinishida qo‘yish, Stanza modelini
        o‘zgartirish, algoritmni shu modelda qo‘llash, to‘g‘riligini 1\,560 masalada tekshirish va vaqt
        yutug‘ini o‘lchash. Algoritm \textbf{nomi} slaydda qoldirildi: komissiyadagi matematiklar uni
        taniydi, nomsiz ko‘rsatilsa, natija o‘zlashtirilgandek ko‘rinishi mumkin. Shuning uchun
        bibliografik iqtiboslar olib tashlandi, nomlar esa qoldi.
  \item \textbf{«Ilk bor» faqat tekshirilsa tasdiqlanadigan joylarda.} UT mualliflari o‘z korpusini
        «birinchi o‘zbek UD korpusi» deb e’lon qilgan (UD v2.15, 2024). Shuning uchun UzUDT uchun
        «badiiy va ta’limiy matnlar uchun ilk bor» va «hozirda eng katta» deyilgan. «O‘zbek tili uchun
        ilk bor MST tahlilchisi» deyilmagan: MST dekodlash umumiy vositalarda keng ishlatiladi.
  \item Vaqt tajribasi: skript --- \texttt{mst\_vaqt\_tajribasi.py},\\ natijalar ---
        \texttt{mst\_vaqt\_natijalari.json}. Seminardan oldin bir marta o‘zingiz ishga tushiring
        (repozitoriy ildizidan, $\approx$1 daqiqa).
  \item Titul slayddagi \textbf{ilmiy rahbar} (prof.~M.\,M.~Aripov deb olindi) va kafedra nomini
        tasdiqlang (\texttt{\textbackslash ilmiyrahbar} buyrug‘i).
  \item \textbf{Muhim.} Maqolalardagi «test» jadvali aslida \textbf{validatsiya} natijalari:
        \texttt{*\_summary.json} dagi \texttt{final\_uas/final\_las} --- dev faylining oxirgi
        qadamdagi qiymati (E2.2 uchun --- eng yaxshi qadam). Taqdimotda raqamlar «validatsiya,
        eng yaxshi nazorat nuqtasi» deb to‘g‘ri belgilangan.
  \item \textbf{Muhim.} Birlashma (.2) tajribalarida parserga \textbf{oltin} UPOS/XPOS/UFeats berilgan
        (\texttt{data/depparse/merged} fayllari), UzUDT (.1) tajribalarida esa eski XLM-R teglovchisining
        bashorati. Shuning uchun maqoladagi «korpuslarni birlashtirish LAS ni 9,6--11,2 punkt
        oshiradi» xulosasi bu jurnallar bilan tasdiqlanmaydi.
  \item .1 va .2 da UzUDT ning \textbf{turli bo‘linmalari}: 435\,/\,48\,/\,198 (eski) va
        451\,/\,45\,/\,188 (yangi); eski dev ning 38 gapi yangi train ichida.
  \item Birlashtirilgan fayllarda kodlash xatosi: «oʻyinni» $\to$ «oÊ»yinni» (train da 2\,554 marta).
  \item E3.2 raqamlari (84,02\,/\,87,07\,/\,70,39\,/\,70,74\,/\,60,05) jurnallarda topilmadi:
        teglovchi jurnali 1\,500-qadamda tugagan, parser jurnali yo‘q.
  \item TahrirchiBERT --- bayt darajasidagi \textbf{BPE} tokenizator (WordPiece emas);
        «bola \#\#lar \#\#ning» misoli haqiqiy emas (haqiqatda \textit{bolalarning} --- bitta token).
  \item Haqiqiy giperparametrlar (W\&B): $\beta_2=0{,}95$, dropout 0,5 (so‘z dropout 0,33),
        bi-affin 400, BiLSTM $3\times400$; maqolada $\beta_2=0{,}999$, dropout 0,33, MLP 500.
        BERT qatlamlari: kodda $(\mathbf H^{(9)}+\mathbf H^{(10)}+\mathbf H^{(11)})/4$.
  \item Bog‘lanish turlari: 39 (maqolada 41). Robustlik jurnalidagi UPOS/UFeats raqamlari
        oltin teglarga nisbatan emas --- taqdimotda faqat UAS/LAS ishlatildi.
\end{enumerate}

\end{spacing}

\newpage

% ------------------------------------------------------------
\section*{1-slayd --- Titul varag‘i}
\timing{0:00}{0:23}

Hurmatli kafedra mudiri! Hurmatli seminar ishtirokchilari!

Sizlarga \say{«O‘zbek tili gaplarini neyron tarmoqlar asosida sintaktik tahlil
qilishning matematik modeli»} mavzusidagi dissertatsiya ishimni taqdim etaman.
Ish matematik modellashtirish, sonli usullar va dasturlar majmui ixtisosligi
bo‘yicha bajarilmoqda.

% ------------------------------------------------------------
\section*{2-slayd --- Korpusdan sintaktik daraxtga}
\timing{0:23}{1:21}

Avval tadqiqot obyektini ko‘rsataman.

\pointto{chap jadval}

Bu --- men yaratgan UzUDT korpusidagi haqiqiy yozuv. Har bir so‘z uchun lemma va
so‘z turkumi berilgan. Eng muhimi --- ikki ustun: \textbf{HEAD} --- so‘zning
\textbf{o‘zagi} raqami va \textbf{DEPREL} --- bog‘lanish turi. Olti yuz sakson to‘rt
gapning har biri qo‘lda shunday belgilangan.

\pointto{daraxt}

O‘ngda --- shu yozuvning grafi. Uchlar --- so‘zlar, yoylar --- bog‘lanishlar.
Ildizdan «beradi» kesimiga yoy chiqadi, undan esa ega, to‘ldiruvchi va boshqa
bo‘laklarga. Har bir so‘zga \textbf{aynan bitta} yoy kiradi, shuning uchun bu ---
daraxt.

Mening vazifam: faqat so‘zlar berilganda shu daraxtni \say{avtomatik topish}.

% ------------------------------------------------------------
\section*{3-slayd --- Daraxt nima uchun kerak: semantika}
\timing{1:21}{2:18}

Nima uchun aynan daraxt? Chunki gapning ma’nosi so‘zlarning o‘zida emas, ular
orasidagi bog‘lanishlarda.

\pointto{1-blok}

Birinchidan, daraxtdan «kim? kimga? nimani?» savollariga javob olinadi: kim ochib
beradi --- asarlar, kimga --- bizga, nimani --- xazina eshigini.

\pointto{2-blok}

Ikkinchidan, so‘zlar bir xil, daraxt boshqa --- \textbf{ma’no boshqa}. «Eski kitob
do‘koni»: birinchi daraxtda kitoblar eski, ikkinchisida --- do‘konning o‘zi eski.

\pointto{3-blok}

Uchinchidan, «Ali kitobni o‘qidi» va «Kitobni Ali o‘qidi»: tartib turli, daraxt bir
xil --- \textbf{ma’no ham bir xil}.

Demak, to‘g‘ri daraxt --- mashina tarjimasi, savol-javob va axborot ajratish
tizimlarining asosi.

% ------------------------------------------------------------
\section*{4-slayd --- Mavzuning dolzarbligi}
\timing{2:18}{2:53}

O‘zbek tili --- \textbf{agglyutinativ} til: «kitoblarimizdagi» so‘zida negiz va to‘rtta
qo‘shimcha bor; eng muhim grammatik ma’lumot so‘z oxirida.

\pointto{diagramma}

Ikkinchi muammo --- ma’lumot kamligi: UD tizimida o‘zbek tili uchun bor-yo‘g‘i
\textbf{bir ming bir yuz sakson to‘rtta} gap bor, turk tilida --- qariyb o‘n ming.
Shuning uchun kichik tanlanmada ishlaydigan, matematik asoslangan model zarur.

% ------------------------------------------------------------
\section*{5-slayd --- Muammoning o‘rganilganlik darajasi}
\timing{2:53}{3:23}

Vaqt chizig‘ida sohaning asosiy bosqichlari ko‘rsatilgan. O‘zbek tili uchun so‘nggi
yillarda TahrirchiBERT modeli va birinchi UD korpusi paydo bo‘ldi.

\pointto{qizil blok}

Biroq korpuslar juda kichik, qism so‘zlardan so‘z vektorini olish operatori
asoslanmagan, dekodlash usullari va modelning turg‘unligi o‘rganilmagan.
\say{Mening ishim shu bo‘shliqni to‘ldiradi.}

% ------------------------------------------------------------
\section*{6-slayd --- Masalaning qo‘yilishi: gapning graf modeli}
\timing{3:23}{4:27}

Endi masalani \textbf{graflar nazariyasi} tilida qo‘yaman.

\pointto{graf}

$n$ ta so‘zli gapga $n+1$ ta uchli to‘liq yo‘naltirilgan graf mos keladi: nolinchi
uch --- ildiz, qolganlari --- so‘zlar. Unda $n$ kvadrat yoy bor, har bir yoyga
\textbf{vazn} beriladi.

Bog‘liqlik daraxti --- shu grafning ildizdan boshlanuvchi \textbf{qamrovchi
daraxti}. Har bir so‘zga yagona yoy kiradi; u boshlangan uch --- shu so‘zning
\textbf{o‘zagi}. Sikl yo‘q, ildizdan bitta yoy chiqadi.

\textbf{Lemma}: bunday daraxtlar soni $n$ ning $n$ minus birinchi darajasiga teng:
o‘n so‘zda --- bir milliard. Men bu formulani $n$ yettigacha barcha daraxtlarni sanab
tekshirdim.

\pointto{qizil blok}

Masala: vaznlar yig‘indisi eng katta bo‘lgan qamrovchi daraxtni topish.

% ------------------------------------------------------------
\section*{7-slayd --- Maqsad, vazifalar, obyekt va predmet}
\timing{4:27}{4:47}

Tadqiqotning maqsadi --- kam resursli sharoitda o‘zbek tili gaplarini tahlil
qilishning \textbf{matematik modeli, sonli usullari va dasturiy majmuasini} yaratish.
Vazifalar slaydda: korpus, graf modeli, modelni o‘zgartirish, sonli usullar va
baholash.

% ------------------------------------------------------------
\section*{8-slayd --- Ilmiy yangilik}
\timing{4:47}{5:38}

Ishning \textbf{to‘rtta} ilmiy yangiligi bor.

\textbf{Birinchi}: o‘zbek tilining badiiy va ta’limiy matnlari uchun \say{ilk bor}
oltin standart UD korpusi --- UzUDT yaratildi.

\textbf{Ikkinchi}: Stanza tahlilchisi o‘zbek tili uchun o‘zgartirildi va \say{ilk bor}
«oxirgi qism so‘z» operatorining ustunligi statistik tasdiqlandi.

\textbf{Uchinchi}: daraxtni izlash murakkabligi $n$ ning $n$ minus birinchi
darajasidan \textbf{$n$ kvadratga} tushirildi; o‘zbek tili uchun \say{ilk bor} graf
va o‘tish usullari bir xil test to‘plamida taqqoslandi.

\textbf{To‘rtinchi}: \say{ilk bor} ikki korpus birlashtirilib o‘qitildi, turg‘unlik
baholandi va dasturiy majmua yaratildi.

% ------------------------------------------------------------
\section*{9-slayd --- Men o‘zgartirgan model}
\timing{5:38}{6:45}

Endi aynan \textbf{qaysi modelni o‘zgartirganimni} aytaman.

Asos --- Stanza kutubxonasining graf asosidagi chuqur bi-affin tahlilchisi. Ko‘k
bloklar --- uning mavjud qismlari, to‘q sariq bloklar --- \textbf{men qo‘shgan} qismlar.

\pointto{to‘q sariq ramka}

Men tahlilchining kirish qatlamiga yangi kanal qo‘shdim: gap TahrirchiBERT
modelidan o‘tadi, qism so‘z vektorlari men kiritgan $\Pi$ operatori orqali so‘z
vektoriga aylanadi. Kirish vektori \textbf{besh yuz} o‘lchamdan \textbf{bir ming bir
yuz qirq uch} o‘lchamga kengaydi; BERT muzlatilgan, o‘qitiladigan parametrlar ---
yigirma yetti million.

\pointto{jadval}

Natijada LAS 55,3 dan 58,2 ga, UPOS 79,2 dan 83,5 ga oshdi.

Bu modelda daraxtni topish uchun barcha variantlarni saralash shart emas ---
\say{$n$ kvadrat amal yetarli}. Buni keyingi slaydlarda ko‘rsataman.

% ------------------------------------------------------------
\section*{10-slayd --- Qism so‘zlarni birlashtirish operatori $\Pi$}
\timing{6:45}{7:24}

BERT so‘zni qism so‘zlarga bo‘ladi, menga esa har bir so‘z uchun bitta vektor
kerak. Buni \textbf{$\Pi$ matritsa} bilan yozdim: so‘z vektori --- qism so‘z
vektorlarining qavariq kombinatsiyasi. Dasturda ikki variantni amalga oshirdim:
o‘rtacha qiymat va «oxirgi qism so‘z».

\pointto{jadval}

TahrirchiBERT bo‘linishida kelishik va shaxs-son qo‘shimchalari doim oxirgi bo‘lakda:
«dagi», «ning», «misiz». Shuning uchun oxirgi bo‘lak grammatik ma’lumotni
yaxshiroq saqlaydi.

% ------------------------------------------------------------
\section*{11-slayd --- Yoy vaznlari: chuqur bi-affin baholash}
\timing{7:24}{8:04}

Yoy vaznlarini neyron tarmoq hisoblaydi. Har bir so‘z ikki rolda tasvirlanadi:
\textbf{tobe} va \textbf{o‘zak}. Juftlik bahosi --- \textbf{bilinear forma}. Ikkinchi
lemmaga ko‘ra bahoning bir qismi softmax da qisqaradi, qolgani --- so‘zning o‘zak
bo‘lishga moyilligi.

\pointto{grafik}

Korpusda yoylarning yarmi qo‘shni so‘zlar orasida, \textbf{sakson to‘rt foizida}
o‘zak uch o‘ngda --- bu o‘zbek tilining SOV tartibi. Natijada grafning vazn
matritsasi $W$ hosil bo‘ladi.

% ------------------------------------------------------------
\section*{12-slayd --- Ikki optimallashtirish masalasi}
\timing{8:04}{8:48}

Seminarda «optimallashtirish nima bo‘yicha?» degan savol berildi. Javob: ishda
\textbf{ikki xil} optimallashtirish masalasi bor.

\pointto{chap ustun}

Birinchisi --- daraxtni topish. Bu diskret masala, u har bir gap uchun yechiladi.
Bu yerda yutuq --- \say{vaqt bo‘yicha}: eksponensial saralash o‘rniga $n$ kvadrat amal.

\pointto{o‘ng ustun}

Ikkinchisi --- parametrlarni topish, ya’ni o‘qitish. Bu uzluksiz masala, u bir marta
yechiladi. Bu yerda yutuq --- \say{aniqlik bo‘yicha}: LAS yigirma to‘rt foizdan ellik
sakkiz foizgacha.

% ------------------------------------------------------------
\section*{13-slayd --- Sonli usul (1): maksimal vaznli qamrovchi daraxt}
\timing{8:48}{9:54}

Birinchi sonli usul. \textbf{Teorema}: model yoylar bo‘yicha faktorlansa, eng ehtimolli
daraxt --- \textbf{maksimal vaznli yo‘naltirilgan qamrovchi daraxt}. Uni Chu--Liu--Edmonds
algoritmi aniq topadi.

\pointto{sxema}

Algoritm: har bir so‘z uchun eng og‘ir kiruvchi yoy tanlanadi. Sikl bo‘lmasa ---
yechim tayyor. Sikl bo‘lsa, u bitta uchga qisqartiriladi va masala qayta yechiladi.

\pointto{«Murakkablik» bloki}

Bitta iteratsiya $n$ kvadrat amal: vazn matritsasining har bir elementi bir marta
ko‘riladi. O‘lchovlarimga ko‘ra iteratsiyalar soni o‘rtacha bir-ikkita. Demak,
murakkablik --- \textbf{$n$ kvadrat}.

\pointto{to‘q sariq blok}

Men bu algoritmni o‘zgartirgan modelimda qo‘lladim va to‘g‘riligini tekshirdim:
\textbf{bir ming besh yuz oltmishta} masalada natija to‘liq saralash bilan \textbf{yuz
foiz} mos keldi.

% ------------------------------------------------------------
\section*{14-slayd --- Hisoblash tajribasi: vaqt bo‘yicha yutuq}
\timing{9:54}{10:50}

Bu yutuq vaqtda qanday ko‘rinadi? Men ikkala usulning vaqtini \textbf{o‘lchadim}.

\pointto{qizil chiziq}

To‘liq saralashda har bir qo‘shimcha so‘z vaqtni taxminan yigirma barobar oshiradi:
sakkiz so‘zli gap --- to‘qqiz soniya, o‘n bir so‘zli o‘rtacha gap --- qariyb
\textbf{ikki sutka}. O‘n to‘qqiz so‘zli gapni saralash esa \textbf{olamning yoshidan}
ham uzoq davom etadi.

\pointto{ko‘k chiziq}

Chu--Liu--Edmonds bilan o‘n bir so‘zli gap --- \textbf{0,42 millisekund}, ya’ni taxminan
350 million marta tez. Butun test to‘plami 0,09 soniyada tahlil qilinadi.

Mana shu --- $n$ ning $n$ minus birinchi darajasidan \say{$n$ kvadratga} o‘tishning
amaliy ma’nosi.

% ------------------------------------------------------------
\section*{15-slayd --- Sonli usul (2): Adam usuli bilan o‘qitish}
\timing{10:50}{11:59}

Ikkinchi sonli usul --- o‘qitish. Noma’lum --- yigirma yetti million parametr,
maqsad funksiyasi --- to‘g‘ri o‘zak va to‘g‘ri bog‘lanish ehtimolliklarining manfiy
logarifmi.

Uni \textbf{Adam} usuli bilan minimallashtirdim: u har bir parametr uchun qadamni
gradientning birinchi va ikkinchi momentlari bo‘yicha moslaydi.

\pointto{to‘q sariq blok}

\textbf{Men bajargan ishlar}: ikki korpus uchun maqsad funksiyasini tuzdim; BERT ni
muzlatib, optimallashtirish o‘lchamini besh barobar kamaytirdim; to‘xtatish qoidasini
qo‘lladim: uch ming qadam yaxshilanish bo‘lmasa --- AMSGrad, yana uch ming qadamdan
keyin --- to‘xtash.

\pointto{grafik}

\say{Natija}: ming qadamda, olti daqiqada, yo‘qotish 14,3 dan 5,4 gacha kamaydi,
LAS esa 24,7 dan \textbf{58,2 foizgacha} oshdi. Keyin LAS tushadi --- bu haddan
tashqari moslashish, shuning uchun eng yaxshi nuqta saqlanadi.

% ------------------------------------------------------------
\section*{16-slayd --- Ma’lumotlar va tajriba dizayni}
\timing{11:59}{12:17}

Tajribalar ikki korpusda o‘tkazildi: men yaratgan \textbf{UzUDT} va \textbf{UT}. Ular
janr jihatidan bir-birini to‘ldiradi. Reja --- uch xil vektor va ikki ma’lumot
sharoitidan iborat faktorli dizayn.

% ------------------------------------------------------------
\section*{17-slayd --- Natijalar}
\timing{12:17}{12:53}

Jadvalda validatsiya natijalari. A blokda TahrirchiBERT va oxirgi qism so‘z
operatori eng yaxshi: LAS --- \textbf{58,23}.

\pointto{bootstrap bloki}

Bu farq tasodifiy emasligini \textbf{juftlashgan bootstrap} bilan tekshirdim: oxirgi
qism so‘z o‘rtacha qiymatdan \textbf{statistik ahamiyatli} yaxshi, $p$ taxminan 0,002.
B blokda ikki korpus birlashmasi; u yerda oltin teglar berilgan, shuning uchun A
bilan bevosita taqqoslanmaydi.

% ------------------------------------------------------------
\section*{18-slayd --- Test to‘plami: dekodlash va turg‘unlik}
\timing{12:53}{13:22}

\pointto{o‘ng jadval}

Bir xil test to‘plamida MST bilan ishlovchi graf modeli ochko‘z o‘tish modelidan
LAS bo‘yicha \textbf{8,7 punkt} yuqori: global optimallashtirishda xatolar zanjir
bo‘ylab tarqalmaydi.

\pointto{grafik}

Turg‘unlik: diakritika yo‘qolsa, LAS atigi bir punkt kamayadi; tinish belgilari
olib tashlansa --- sakkiz punkt.

% ------------------------------------------------------------
\section*{19-slayd --- Dasturiy majmua}
\timing{13:22}{13:44}

Barcha usullar dasturiy majmuada amalga oshirildi: Stanza'ga kiritgan
o‘zgarishlarim, spaCy uchun o‘zbek tili moduli, baholash va bootstrap vositalari ---
jami \textbf{uch ming qatordan ortiq} Python kodi. Korpus UD da, modellar HuggingFace
da ochiq.

% ------------------------------------------------------------
\section*{20-slayd --- Xulosa}
\timing{13:44}{14:38}

Xulosa qilib aytganda:

\begin{itemize}
  \item badiiy va ta’limiy matnlar uchun \textbf{ilk bor} UzUDT korpusini yaratdim;
  \item sintaktik tahlilni \textbf{graflar nazariyasi} tilida --- maksimal vaznli qamrovchi
        daraxt masalasi sifatida qo‘ydim;
  \item Stanza tahlilchisini o‘zgartirdim; oxirgi qism so‘z operatori LAS ni 2,6 punkt
        oshiradi;
  \item o‘zgartirgan modelimda daraxtni izlash \textbf{$n$ kvadrat} amal bilan bajariladi:
        o‘rtacha gap uchun ikki sutka o‘rniga yarim millisekund;
  \item Adam bilan o‘qitishda ming qadamda LAS 58 foizga yetdi;
  \item dasturiy majmua yaratildi.
\end{itemize}

\pausehere

Kelgusida barcha modellarni \textbf{yagona test to‘plamida} qayta baholash
rejalashtirilgan.

% ------------------------------------------------------------
\section*{21-slayd --- Aprobatsiya va nashrlar}
\timing{14:38}{14:53}

Natijalar oltita ilmiy ishda e’lon qilingan, barchasi Scopus bazasida, jumladan
\textit{Data in Brief} jurnalida (Q1) va Springer NLDB to‘plamida (Q2).

% ------------------------------------------------------------
\section*{22-slayd --- E’tiboringiz uchun rahmat!}
\timing{14:53}{14:59}

\say{E’tiboringiz uchun rahmat! Savollaringizga javob berishga tayyorman.}

{\small\itshape Zaxira slayd (23-bet, BiLSTM formulalari) --- faqat kodlovchi haqida savol
berilsa ko‘rsatiladi.}

% ============================================================
\newpage
\section*{Ehtimoliy savollar va qisqa javoblar}

\begin{itemize}
  \item \textit{«O‘zak» deganda nimani tushunasiz?} --- Yoy $h\to i$ ning boshlang‘ich uchi:
        $i$-so‘zga kiruvchi yagona yoy shu uchdan chiqadi. Tilshunoslikdagi «hokim» so‘ziga mos.
        Morfologik o‘zak bilan adashtirmaslik uchun 4-slayddagi morfema sxemasida «negiz» deyildi.
  \item \textit{$n^{n-1}$ dan $n^2$ ga o‘tish --- sizning natijangizmi? Algoritm kimniki?} ---
        Algoritmning o‘zi klassik: Chu--Liu--Edmonds algoritmi. Mening hissam: o‘zbek tili
        gaplarini tahlil qilish masalasini aynan shu graf masalasi ko‘rinishida qo‘ydim, uni o‘zim
        o‘zgartirgan modelda qo‘lladim, to‘g‘riligini 1\,560 masalada to‘liq saralash bilan
        tekshirdim va vaqt yutug‘ini o‘lchadim: o‘rtacha gap uchun $\approx1{,}7$ sutka o‘rniga 0,42 ms.
  \item \textit{Optimallashtirish nima bo‘yicha?} --- Ikki xil: (P1) daraxtni topish --- vaqt
        (amallar soni) bo‘yicha: $n^{n-1}$ variant o‘rniga $O(n^2)$ amal; (P2) o‘qitish --- yo‘qotish
        funksiyasi, ya’ni aniqlik bo‘yicha: LAS $24{,}7\to58{,}2\,\%$.
  \item \textit{CLE ning eng yomon holatdagi murakkabligi qancha?} --- Stanza amalga oshirishida
        bitta iteratsiya $O(n^2)$; eng yomon holatda iteratsiyalar soni $n$ gacha, ya’ni $O(n^3)$;
        bitta ildiz sharti uchun nomzod ildizlar bo‘yicha qayta yechish qo‘shiladi. Amalda (o‘lchov)
        iteratsiyalar 1--2, vaqt $\propto n^{1,56}$. Zich graflar uchun nazariy $O(n^2)$ amalga
        oshirish ham mavjud (Tarjan, 1977).
  \item \textit{Butun modelning murakkabligi qanday?} --- BERT: $O(m^2)$ ($m$ --- qism so‘zlar soni),
        BiLSTM: $O(n)$, bi-affin baholash: $O(n^2)$, MST: amalda $O(n^2)$ $\Rightarrow$ umumiy $O(n^2)$.
  \item \textit{Vaqtni tasodifiy vaznlarda o‘lchash to‘g‘rimi?} --- Tasodifiy vaznlarda ochko‘z qadam
        ko‘p sikl hosil qiladi, ya’ni bu og‘ir holat. O‘qitilgan model vaznlari aniqroq, sikllar
        kamroq --- vaqt bundan ham kichik. To‘liq saralash vaqti esa vaznlarga bog‘liq emas, faqat $n$ ga.
  \item \textit{Adam ni o‘zgartirdingizmi? Unda siz nima qildingiz?} --- Adam usulining o‘zi standart.
        Men: (1) ikki korpus uchun maqsad funksiyasini tuzdim; (2) BERT ni muzlatib, optimallashtirish
        o‘lchamini 137 mln dan 27,6 mln ga ($\approx$5 marta) kamaytirdim; (3) validatsiyaga asoslangan
        to‘xtatish qoidasini qo‘lladim (3\,000 qadam $\to$ AMSGrad $\to$ 3\,000 qadam $\to$ to‘xtash);
        (4) o‘qitishni jurnallashtirib, 5 tajribada yaqinlashishni tahlil qildim: eng yaxshi nuqta
        800--1\,300 qadamda.
  \item \textit{Nima uchun Adam, SGD emas?} --- So‘z va lemma vektorlari jadvalida gradientlar siyrak:
        kam uchraydigan so‘z shakllari kam yangilanadi. Adam har bir koordinata uchun qadamni
        $\sqrt{\hat v}$ ga bo‘ladi, shuning uchun kam yangilanadigan parametrlar kattaroq effektiv qadam
        oladi. Agglyutinativ tilda bu muhim.
  \item \textit{Nima uchun $\beta_2=0{,}95$, standart 0,999 emas?} --- Ikkinchi momentning «xotirasi»
        taxminan $1/(1-\beta_2)$ qadam: 0,95 da --- 20 qadam, 0,999 da --- 1\,000 qadam. Qisqa xotira
        qadam hajmini gradient o‘zgarishiga tezroq moslaydi; bu bi-affin tahlilchilar uchun tavsiya
        etilgan qiymat (Stanza sozlamasi).
  \item \textit{Bir qadamda qancha ma’lumot ishlatiladi?} --- Partiya 5\,000 so‘zgacha. UzUDT o‘quv
        to‘plami 435 gap ($\approx$5\,700 so‘z, ildiz va qo‘shimcha namunalar bilan) --- 1 davr = 2 qadam;
        birlashmada 3 qadam. Demak, 1\,000 qadam $\approx$ 500 davr.
  \item \textit{AMSGrad nima beradi?} --- Maxrajda $v_t$ o‘rniga $\max_{s\le t}v_s$ ishlatiladi, shuning
        uchun effektiv qadam o‘smaydi; bu Adam ning yaqinlashmaslik holatlarini bartaraf etadi. Bizda u
        eng yaxshi nuqtadan keyin, yakuniy bosqichda ishlaydi.
  \item \textit{Nega har bir so‘z uchun alohida softmax (lokal normallash), global ehtimollik
        modeli emas?} --- Lokal model o‘qitishni sodda va tez qiladi. Dekodlash baribir global: MST
        daraxt shartini kafolatlaydi. Global model uchun bo‘linish funksiyasi $Z(x)$ Kirchhoff
        matritsa-daraxt teoremasi bilan determinant orqali $O(n^3)$ da hisoblanadi --- bu keyingi bosqich.
  \item \textit{CLE algoritmi nega optimal?} --- Asosiy lemma: ochko‘z tanlovdagi sikl $C$ optimal
        daraxtda faqat bitta kiruvchi yoy bilan almashadi; qisqartirilgan grafdagi optimal daraxt asl
        grafdagi optimal daraxtga mos keladi. Bundan tashqari, men buni $n\le7$ da 1\,560 masalada
        to‘liq saralash bilan solishtirib tekshirdim.
  \item \textit{Bitta ildiz sharti qanday hisobga olinadi?} --- Agar ildizdan bir nechta yoy chiqsa,
        nomzod ildizlar bo‘yicha CLE qayta yechiladi va eng katta vaznli daraxt tanlanadi.
  \item \textit{Nima uchun Koshi yadrosi, Gauss emas?} --- Yoy uzunligi taqsimoti og‘ir dumli
        (3,3\,\% yoylar 10 dan uzun). Gauss yadrosi uzun yoylarni haddan tashqari jazolaydi, Koshi
        yadrosi esa $1/d^2$ tezlikda so‘nadi.
  \item \textit{Nima uchun $\tfrac14$ koeffitsient uchta qatlam uchun?} --- Bu Stanza kutubxonasining
        standart sozlamasi: 9--11-qatlamlar yig‘indisi 4 ga bo‘linadi. O‘zgarmas masshtab keyingi
        chiziqli qatlamda to‘liq kompensatsiya qilinadi, ya’ni model sinfi o‘zgarmaydi.
  \item \textit{Nega BERT sozlanmagan (muzlatilgan)?} --- 781 gapda 110 mln parametrni sozlash haddan
        tashqari moslashish xavfini oshiradi. Muzlatilganda o‘qitiladigan parametrlar 27,6 mln,
        o‘qitish barqaror va tez (eng yaxshi model 6 daqiqada).
  \item \textit{Natijalar nega validatsiya to‘plamida?} --- Repozitoriy jurnallarida o‘qitish davomida
        validatsiya ko‘rsatkichlari saqlangan. Test to‘plamidagi baholash E2.1 va spaCy uchun bajarilgan
        (18-slayd); barcha modellarni yagona test to‘plamida qayta baholash --- eng yaqin vazifa.
  \item \textit{Statistik ahamiyatlilik qanday baholandi?} --- Juftlashgan bootstrap: birlik --- gap
        (xatolar gap ichida bog‘liq), $B=10^4$, bir tomonlama $p=P^{*}(\Delta^{*}\le0)$.
  \item \textit{Korpuslarni birlashtirish qancha foyda beradi?} --- Birlashma tajribalarida parserga
        oltin teglar berilgan va validatsiya to‘plamlari ham farqli. Shuning uchun sof ta’sirni
        ajratish uchun qo‘shimcha tajriba (bir xil teglar, umumiy test gaplari) rejalashtirilgan.
  \item \textit{«Ilk bor» deyishga asos bormi? UT ham «birinchi o‘zbek UD korpusi»-ku.} --- Ha, UT ---
        birinchi o‘zbek UD korpusi (2024, yangilik va badiiy matnlar). Shuning uchun men «badiiy va
        ta’limiy matnlar uchun ilk bor» deyman; UzUDT --- hozirda eng katta ochiq o‘zbek UD korpusi.
        Ikki korpusni birlashtirib o‘qitish va graf/o‘tish usullarini bir xil test to‘plamida taqqoslash
        esa UzUDT e’lon qilingandan keyingina mumkin bo‘ldi --- buni birinchi bo‘lib men bajardim.
  \item \textit{Nima uchun o‘tish modeli (spaCy) UPOS da yaxshi, LAS da yomon?} --- spaCy transformerni
        teglash bilan birga sozlaydi, shuning uchun teglash yaxshi. Lekin ochko‘z ketma-ket qarorlar
        xatoni tarqatadi va global optimal yechim kafolatlanmaydi.
  \item \textit{Katta til modellari (GPT va boshqalar) bilan taqqoslanganmi?} --- Hozircha yo‘q; bu
        kelgusi ish. Til modelining javobida daraxt sharti kafolatlanmaydi, bizning yondashuvda esa MST
        orqali kafolatlanadi.
  \item \textit{Mavzuda nega «semantik» tahlil yo‘q?} --- Tadqiqot morfologik va sintaktik qatlamlarga
        qaratilgan. Sintaktik daraxt --- semantik tahlil uchun asos (3-slayddagi misollar); semantik
        rollar annotatsiyasi --- keyingi bosqich.
  \item \textit{Amaliy qo‘llanilishi?} --- Mashina tarjimasi, axborot qidirish, savol-javob tizimlari,
        grammatik tekshiruv va ta’lim matnlarini tahlil qilish.
\end{itemize}

\end{document}
